\documentclass[10pt,a4paper]{amsart}
\usepackage[margin=19mm]{geometry}
\usepackage{amsmath,amssymb,mathtools}
\usepackage[scr=rsfs,cal=euler]{mathalfa}
\usepackage{xfrac}
\usepackage[unicode,hidelinks,bookmarks=false]{hyperref}
\newcommand{\ie}{i.e.\ }
\newcommand{\cf}{cf.\ }
\newcommand{\resp}{respectively\ }
\newcommand{\Subsection}{\subsection}
\providecommand{\nobreakdash}{\nobreak\hbox{-}}
\setlength{\emergencystretch}{2em}
\multlinegap0pt
\usepackage{amssymb}
\usepackage[shortlabels]{enumitem}
\usepackage{xcolor}
\newtheorem{lemma}{Lemma}
\newtheorem{proposition}{Proposition}
\DeclareMathOperator{\Sat}{Sat}
\newcommand{\symd}{\mathbin{\triangle}}
\title{Degreewise Cut-Semigroup Saturation for $K_5$-Minor-Free Graphs and Seymour's Planar Edge-Colouring Conjecture}
\author{SeungJu Lee}
\date{Source: arXiv:2609.12732v1 (CC BY 4.0)}
\begin{document}
\maketitle

\noindent\emph{Source and licence.} Degreewise Cut-Semigroup Saturation for $K_5$-Minor-Free Graphs and Seymour's Planar Edge-Colouring Conjecture. Version arXiv:2609.12732v1 (CC BY 4.0), distributed by arXiv under the Creative Commons Attribution 4.0 International licence, http://creativecommons.org/licenses/by/4.0/, as recorded in the arXiv licence metadata for this exact version and shown on the abstract page https://arxiv.org/abs/2609.12732v1. Full licence notice: \texttt{licenses/arxiv-2609.12732-CC-BY-4.0.txt}.\par\smallskip

\noindent\textit{Editorial scope.} Counted unit: the article's proposition prop:reverse (source0.tex lines 407--409). The article's complete original proof of it is delivered with it (source0.tex lines 410--432, one \texttt{proof} environment of 23 lines). No mathematics is written, repaired or re-derived: the statement and the proof are the article's own lines, in the article's own notation, and no retained line is substituted or re-flowed.  Retained uncounted as the source-local context the proof needs: lines 77--79; lines 369--371; lines 395--398.  Nothing was omitted from any retained span, so the package carries no omission record.  No retained line contains a citation, so the wrapper carries no bibliography and no bibliography entry.  No retained line contains a figure environment, an image-inclusion command or drawing code, so no image is delivered and the package declares no asset dependency.  Editorial formatting only, in addition: an AMS \texttt{amsart} preamble that restates the article's own declarations and macros used by the retained lines, namely the environments \texttt{lemma}, \texttt{proposition} on separate counters, together with the standard packages those lines need.  The retained environments print in this document as Lemma 1, Proposition 1 in order of appearance, whereas the article prints the same items with the numbering of its own full document.  The complete native source of the article as distributed in the arXiv e-print for this exact version is bundled unchanged as \texttt{sources/210146/source0.tex}.
\begin{equation}\label{eq:oddcut}
 |\delta_R(X)|\geq k\quad\text{for every }X\subseteq V(R)\text{ of odd cardinality}.
\end{equation}

\begin{equation}\label{eq:switch}
 p_e=\begin{cases}k-m_e,&e\in M_0,\\m_e,&e\notin M_0.\end{cases}
\end{equation}

\begin{lemma}\label{lem:simplify}
The vector $\bar p$ belongs to $\Sat_k(G)$.
Every decomposition of $\bar p$ into $k$ cuts of $G$ lifts to a decomposition of $p$ into $k$ Eulerian edge sets of $H$.
\end{lemma}

\begin{proposition}\label{prop:reverse}
If every simple planar graph is saturated at height $k$, then every planar $k$-graph is $k$-edge-colourable.
\end{proposition}

\begin{proof}
First consider connected nonempty $R$.
By the assumption and Lemma~\ref{lem:simplify}, write $p=\sum_{i=1}^k\chi_{C_i}$ with each $C_i$ Eulerian in $H$.
Define $Q_i=M_0\symd C_i$.
Each $Q_i$ has odd degree at every vertex.
For every edge, \eqref{eq:switch} yields
\begin{equation}\label{eq:aggregate}
 \sum_{i=1}^k\chi_{Q_i}(e)=m_e.
\end{equation}
For $e\in M_0$, the left side is $k-p_e$; otherwise it is $p_e$.
Consequently,
\begin{equation}\label{eq:budget}
 \sum_{i=1}^k d_{Q_i}(v)=m(\delta_H(v))=k\quad(v\in V(H)).
\end{equation}
The left side consists of $k$ positive odd integers.
All are therefore one, so every $Q_i$ is a perfect matching.
For each support edge, assign its parallel copies bijectively to the indices where it occurs in \eqref{eq:aggregate}.
The resulting $k$ perfect matchings partition $E(R)$.

Every component of a disconnected $k$-graph has even order: an odd component would violate \eqref{eq:oddcut}.
Each component inherits the same condition, so the connected construction applies componentwise with the same colour set.
The empty graph is immediate.
\end{proof}

\end{document}
